\subsubsection{Quantitative Evaluierung}

Um über die Segmentierungsgüte über die qualitative Sichtprüfung hinaus objektive zu bewerten, schließt sich eine quantitative Evaluierung an. Das implementierte Verfahren gleicht die Position des städtischen Baumkatasters als Refernzdaten mit den detektierten Baumspitzen ab. Die algorithmische Umsetzung dieses Abgleichs entspricht dem in Kapitel 4.3.2 beschrieben Vorgang zu Datenanreicherung, inklusive euklidischer 3D-Distanz und globaler Zuordnung über die ungarische Methode. Ein direkter Abgleich der Daten wird jedoch durch die Unvollständigkeit des Referenzdatensatzes erschwert. Denn das Baumkataster verzeichnet ausschließlich Bäume im öffentlichen Raum, während die Überfliegung mit dem LiDAR-Sensor auch Baumpunkte auf Privat- oder Gewerbeflächen aufgenommen werden. Um diese Punkte nicht als Fehldetektionen zu werten, wird vor dem Zuordungsprozess eine räumliche Filterung auf die LiDAR-Punktwolke angewandt. Um die Referenzbäume wird ein Puffer von \textit{8 m} gelegt, die zu einer gemeinsamen validen Fläche verschmelzen. Für die statistische Auswertung werden ausschließlich Baumspitzen innerhalb dieser Maske berücksichtigt \cite{munzingerMappingUrbanForest2022}.

Die verbleibenden Paare innerhalb der Validierungszone werden für die Berechnung standardisierter Gütemaße in folgende Kategorien unterteilt und als \textit{JSON}-Datei gespeichert.
\begin{itemize}
    \item \textit{True Positives (TP)} entspricht der Anzahl korrekt detektierte Baumspitzen, die innerhalb des Suchradius einem Katasterbaum zugeordnet wurden.
    \item \textit{False Positives (FP)} zählt fälschlicherweise detektierte Baumspitzen ohne korrespondierenden Katastereintrag. Hohe Werte treten primär als Folge von Übersegmentierung auf.
    \item \textit{False Negatives (FN)} repräsentieren nicht erfasste Katasterbäume, was meist ein Resultat von Untersegmentierung ist.
\end{itemize}
Aus diesen absoluten Werten können Kennzahlen abgeleitet werden, welche die Ergebnisse normieren und die statistische Leistungsfähigkeit des Workflows als Verhältniswerte beschreiben. Die Genauigkeit (\textit{Precision}) bestimmt das Verhältnis der korrekt detektierten Bäume zur Gesamtzahl aller gefunden Objekte und wird bestimmt mittel $\frac{TP}{TP + FP}$. Eine hohe \textit{Precision} zeigt, dass der Workflow kaum Fehlinterpretation erzeugt. Die Trefferquote (\textit{Recall}) setzt die korrekt detektierten Bäume in Relation zu allen vorhandenen Katasterbäumen und belegt den Grad der Gesamterfassung mit der Formel $\frac{TP}{TP + FN}$. Der \textit{F1-Score} bildet das harmonische Mittel aus \textit{Precision} und \textit{Recall}. Er dient als globaler Qualitätsindikator, der die Balance zwischen den beiden Metriken bewertet. Ein hoher \textit{F1-Score} wird nur dann erreicht, wenn sowohl eine hohe Trefferquote als auch eine minimale Fehlinterpretationen erzielt wird \cite{zhangIndividualTreeSegmentation2024}.
